\documentclass[10pt, aspectratio=169]{beamer}

% ── Paquetes ────────────────────────────────────────────────────────────────
\usepackage[utf8]{inputenc}
\usepackage[T1]{fontenc}
\usepackage[spanish]{babel}
\usepackage{graphicx}
\usepackage{booktabs}

% Colores
\definecolor{headercolor}{RGB}{45, 62, 80}
\definecolor{rowcolor1}{RGB}{245, 247, 250}
% ── Tema de Beamer ──────────────────────────────────────────────────────────
\usetheme{Madrid}
\usecolortheme{default}
\definecolor{azulUVa}{RGB}{0, 51, 153}
\definecolor{rojoUVa}{RGB}{193, 0, 60}
\setbeamercolor{palette primary}{bg=azulUVa,fg=white}
\setbeamercolor{palette secondary}{bg=rojoUVa,fg=white}
\setbeamercolor{palette tertiary}{bg=black,fg=white}
\setbeamercolor{structure}{fg=azulUVa} % Itemize, enumerate, etc

% ── Datos de la Portada ─────────────────────────────────────────────────────
\title[Práctica 2: Locust y Prometheus]{Práctica 2: Pruebas de carga, Locust y Prometheus}
\subtitle{Evaluación de Sistemas Informáticos}
\author[Grupo 16]{Grupo 16}
\date{\today}

% ════════════════════════════════════════════════════════════════════════════
\begin{document}
	
	% ── Diapositiva 1: Portada ──────────────────────────────────────────────────
	\begin{frame}
		\titlepage
	\end{frame}
	
	% ── Diapositiva 2: Índice ───────────────────────────────────────────────────
	\begin{frame}{Índice}
		\tableofcontents
	\end{frame}
	
	% ════════════════════════════════════════════════════════════════════════════
	\section{Descripción general del entorno}
	
	\begin{frame}{¿Qué elementos son clave?}
		\begin{itemize}
			\item Servidor de autenticación OAuth2 desarrollado en Python (Flask y SQLite)
			\item Implementa flujos clave del RFC 6749: \textit{Authorization Code}, \textit{Client Credentials} y \textit{Refresh Token}.
			\item Ejecutado mediante el gestor \texttt{uv} en el puerto 5001.
		\end{itemize}
	\end{frame}
	
	\section{¿Cómo hemos trabajado?}
	
	\begin{frame}{Distribución del trabajo}
		\centering
		\renewcommand{\arraystretch}{1.3}
		
		\begin{tabular}{p{3.5cm} p{2.5cm} p{5cm}}
			\toprule
			
			\textcolor{white}{\textbf{Nombre}} & 
			\textcolor{white}{\textbf{\% Contribución}} & 
			\textcolor{white}{\textbf{Responsabilidades}} \\
			
			\midrule
			
			Marcos & 33\% & Desarrollo de scripts del fichero Locust y redacción del cuaderno de bitácora \\
			
			Edgar & 33\% & Desarrollo de las metricas y de \textit{models.py},conexión con el Grafana y documentación \\
			
			Álvaro & 33\% & Gestión de la VM y del repositorio de \hyperlink{https://github.com/ALVPI/ESI/tree/code}{GitHub}, modificación de \textit{oauth2\_routes}, creación del \textit{test.py} y  documentación \\
			
			\bottomrule
		\end{tabular}
		
	\end{frame}
	% ════════════════════════════════════════════════════════════════════════════
	\section{Modificaciones en el código para la práctica}
	
	\begin{frame}{Modificaciones en el Backend}
		La instrumentación se ha integrado siguiendo el principio DRY (intentando evitar los \textit{circular imports}):
		\vspace{0.3cm}
		\begin{itemize}
			\item \textbf{\texttt{metrics.py}:} Centralización de métricas (Counters, Gauges, Histograms) para evitar importaciones circulares. Uso estratégico de etiquetas (ej. \texttt{grant\_type})
			\item \textbf{\texttt{models.py}:} Auditoría de la base de datos (códigos generados/consumidos, detección de tokens revocados).
			\item \textbf{\texttt{oauth2\_routes.py}:} Seguimiento del ciclo de vida de los flujos de autenticación (iniciados, exitosos, fallidos).
		\end{itemize}
	\end{frame}
	
	% ════════════════════════════════════════════════════════════════════════════
	\section{Enfoque del monitoreo}
	
	\begin{frame}{Arquitectura de Monitorización}
		\textbf{Enfoque EDR (\textit{Endpoint Detection and Response}):}
		\begin{itemize}
			\item \textbf{Rendimiento:} Histogramas para evaluar tiempos de respuesta y control de peticiones simultáneas.
			\item \textbf{Estado de Tokens:} Monitorización de tokens activos, revocados, códigos pendientes y expirados.
			\item \textbf{Seguridad:} Detección de ataques de fuerza bruta (M2M y Web) e intentos de \textit{Token Replay}.
		\end{itemize}
	\end{frame}
	
	% ════════════════════════════════════════════════════════════════════════════
	\section{Pruebas de Carga con Locust}
	
	\begin{frame}{Diseño de Escenarios en Locust}
		En el \texttt{locustfile.py} hemos intentado simular un comportamiento "real"  \textbf{50\% tráfico legítimo, 50\% tráfico malicioso}.
		
		\vspace{0.3cm}
		\begin{table}[H]
			\centering
			\small
			\begin{tabular}{llcc}
				\toprule
				\textbf{Tarea} & \textbf{Tipo} & \textbf{Peso} & \textbf{Prob.} \\
				\midrule
				\texttt{valid\_client\_credentials} & Tráfico legítimo M2M    & 4 & 40\% \\
				\texttt{brute\_force\_api}          & Ataque fuerza bruta M2M & 2 & 20\% \\
				\texttt{brute\_force\_user\_login}  & Ataque fuerza bruta web & 2 & 20\% \\
				\texttt{valid\_user\_login}         & Tráfico legítimo web    & 1 & 10\% \\
				\texttt{replay\_attack}             & Ataque token replay     & 1 & 10\% \\
				\bottomrule
			\end{tabular}
		\end{table}
	\end{frame}
	
	\begin{frame}{Ejecución y Verificación}
		\textbf{Configuración de Pruebas:}
		\begin{itemize}
			\item \textbf{Carga normal:} 50 usuarios, \textit{spawn rate} 5/s durante 5 min].
			\item \textbf{Prueba de estrés:} 200 usuarios, \textit{spawn rate} 10/s durante 3 min].
		\end{itemize}
		
		\vspace{0.5cm}
		\textbf{Garantías del sistema:}
		\begin{enumerate}
			\item Uso de script de verificación (\texttt{test\_metrics.py}) para validar manualmente los 9 escenarios de métricas antes del ataque global.
			\item Limitación de almacenamiento en Prometheus (\texttt{retention.size=500MB}) para proteger los recursos de la máquina virtual.
		\end{enumerate}
	\end{frame}
\section{Grafana}
	\begin{frame}
		\centering
		\Huge \textbf{\hyperlink{http://virtual.lab.inf.uva.es:31162/d/ad46rtj/esi-p2-16?orgId=1&from=now-5m&to=now&timezone=browser}{Nuestro panel de Grafana}}\\
		\vspace{1cm}
	\end{frame}
	
	% ── Diapositiva Final ───────────────────────────────────────────────────────
	\begin{frame}
		\centering
		\Huge \textcolor{azulUVa}{\textbf{¡Gracias por su atención!}}\\
		\vspace{1cm}
		\Large ¿Preguntas?
	\end{frame}
	
\end{document}